% 6.2 双点群的对称适配函数
\section{双点群的对称适配函数}

第二章我们考虑了点群对称适配函数的推导；这些函数已在表 2.4--2.6 中显式列出。这样的函数或者是一个单独的球谐函数 $Y^{m}_{l}(\theta, \phi)$，或者是具有固定 $l$ 值、不同 $m$ 值的球谐函数的某个线性组合。在忽略自旋的近似下，这些对称适配函数对分子与固体的计算有用。若想恰当地计入自旋，如本章前面已经指出的，就必须考虑点群的双值表示。确定这些双值表示的基函数的问题已由多位作者考虑过（Cracknell 1969\textit{a}, Cracknell and Joshua 1970, Onodera and Okazaki 1966, Teleman and Glodeanu 1967）。

式 (2.1.1) 定义的球谐函数 $Y^{m}_{l}(\theta, \phi)$ 是量子力学算符 $L^{2}$（轨道角动量平方）与 $L_{z}$（角动量 $z$ 分量）的本征函数，对应于本征值 $l(l+1)\hbar^{2}$ 与 $m\hbar$。为了把自旋包括进来，需要把 $L^{2}$ 与 $L_{z}$ 换成总角动量算符 $J^{2}$ 与 $J_{z} = L_{z} + S_{z}$；于是需要用球谐旋量（spherical spinor）$\langle l\,m;\,s\,m_{s} |$ 或形如 $\sum c\,Y^{m}_{l}\chi^{m_{s}}_{s}$ 的线性组合作为基。与第二章的做法完全平行，把算符 $\mathbf{P}$、$\mathbf{Q}$、$\mathbf{R}$、$\mathbf{S}$ 的矩阵代表（表 5.1 的钥匙）作用在这些基上，即可生成双值表示的对称适配线性组合。表 6.5 与 6.6 按点群逐一列出了这些函数（原书扫描见下）。对双点群对称适配函数的详细处理与表，见下述扫描页。

\begin{figure}[H]
\centering
\includegraphics[width=0.86\textwidth]{c6_62_p450.jpg}
\caption*{\textbf{表 6.5}\ 双值点群的对称适配函数（原书第 437 页起扫描）}
\end{figure}
\begin{figure}[H]
\centering
\includegraphics[width=0.86\textwidth]{c6_62_p451.jpg}
\caption*{原书第 438 页}
\end{figure}
\begin{figure}[H]
\centering
\includegraphics[width=0.86\textwidth]{c6_62_p452.jpg}
\caption*{原书第 439 页}
\end{figure}
\begin{figure}[H]
\centering
\includegraphics[width=0.86\textwidth]{c6_62_p453.jpg}
\caption*{原书第 440 页}
\end{figure}
\begin{figure}[H]
\centering
\includegraphics[width=0.86\textwidth]{c6_62_p454.jpg}
\caption*{原书第 441 页}
\end{figure}
\begin{figure}[H]
\centering
\includegraphics[width=0.86\textwidth]{c6_62_p455.jpg}
\caption*{原书第 442 页}
\end{figure}
\begin{figure}[H]
\centering
\includegraphics[width=0.86\textwidth]{c6_62_p456.jpg}
\caption*{原书第 443 页}
\end{figure}
\begin{figure}[H]
\centering
\includegraphics[width=0.86\textwidth]{c6_62_p457.jpg}
\caption*{原书第 444 页}
\end{figure}
\begin{figure}[H]
\centering
\includegraphics[width=0.86\textwidth]{c6_62_p458.jpg}
\caption*{原书第 445 页}
\end{figure}
\begin{figure}[H]
\centering
\includegraphics[width=0.86\textwidth]{c6_62_p459.jpg}
\caption*{原书第 446 页}
\end{figure}
\begin{figure}[H]
\centering
\includegraphics[width=0.86\textwidth]{c6_62_p460.jpg}
\caption*{原书第 447 页}
\end{figure}
\begin{figure}[H]
\centering
\includegraphics[width=0.86\textwidth]{c6_62_p461.jpg}
\caption*{原书第 448 页}
\end{figure}
\begin{figure}[H]
\centering
\includegraphics[width=0.86\textwidth]{c6_62_p462.jpg}
\caption*{原书第 449 页}
\end{figure}
\begin{figure}[H]
\centering
\includegraphics[width=0.86\textwidth]{c6_62_p463.jpg}
\caption*{原书第 450 页}
\end{figure}
\begin{figure}[H]
\centering
\includegraphics[width=0.86\textwidth]{c6_62_p464.jpg}
\caption*{原书第 451 页}
\end{figure}
\begin{figure}[H]
\centering
\includegraphics[width=0.86\textwidth]{c6_62_p465.jpg}
\caption*{原书第 452 页}
\end{figure}
\begin{figure}[H]
\centering
\includegraphics[width=0.86\textwidth]{c6_62_p466.jpg}
\caption*{原书第 453 页}
\end{figure}

% 6.3 空间群的双值表示
\section{空间群的双值表示}

我们在前两节中已经看到：在具有某个点群 $\mathbf{G}$ 对称性的系统中，半奇整数自旋的实体需要使用双群 $\mathbf{G}^{\dagger}$ 或 $\mathbf{G}$ 的双值表示。类似地，如果在具有 230 个空间群之一的对称性的系统中研究半奇整数自旋的实体，就必须研究空间群的双值表示。因此，当在确定晶体固体电子能带结构所用的哈密顿量中计入依赖自旋的项时，双值空间群表示是重要的（Elliott 1954\textit{e}）。必须把定义 6.1.1 的范围扩展到空间群。空间群 $\mathbf{G}$ 由元素 $\{R \mid \mathbf{v}\}$ 组成，它们满足乘法规则
\begin{equation*}
\{R_{2} \mid \mathbf{v}_{2}\}\{R_{1} \mid \mathbf{v}_{1}\} = \{R_{2}R_{1} \mid \mathbf{v}_{2} + R_{2}\mathbf{v}_{1}\}
\tag{1.5.3}
\end{equation*}
并构成群；$\mathbf{G}$ 可以用某个布拉维格子平移群 $\mathbf{T}$ 的左陪集代表表示为
\begin{equation*}
\mathbf{G} = \{R_{1} \mid \mathbf{v}_{1}\}\mathbf{T} + \{R_{2} \mid \mathbf{v}_{2}\}\mathbf{T} + \dots + \{R_{h} \mid \mathbf{v}_{h}\}\mathbf{T}.
\tag{3.5.2}
\end{equation*}
转动部分 $R_{1}, R_{2}, \dots, R_{h}$ 构成 32 个晶体学点群之一，构造其双群的规则由定义 6.1.1 给出。与该点群的每个元素 $R_{i}$ 相对应，双点群中有两个元素，比如说 $R_{i}$ 与 $\overline{R}_{i}$。考虑双空间群时，$R_{i}$ 与 $\overline{R}_{i}$ 都被看作对矢量 $\mathbf{v}_{j}$ 有同样的作用，即
\begin{equation*}
\overline{R}_{i}\mathbf{v}_{j} = R_{i}\mathbf{v}_{j}.
\tag{6.3.1}
\end{equation*}
在构成群乘积时，$R_{i}$ 与 $\overline{R}_{i}$ 按双点群的乘法规则相乘，也就是说，在与 $SU(2)$ 子群的同构中 $R_{i}$ 与 $\overline{R}_{i}$ 分别对应于 $+\mathbf{u}$ 与 $-\mathbf{u}$。于是可以定义双空间群如下：

\noindent\textbf{定义 6.3.1}\quad 由式 (3.5.2) 定义的空间群 $\mathbf{G}$ 的双群 $\mathbf{G}^{\dagger}$ 由下式给出：
\begin{equation*}
\begin{array}{rl}
\mathbf{G}^{\dagger} &= \{R_{1} \mid \mathbf{v}_{1}\}\mathbf{T} + \{\overline{R}_{1} \mid \mathbf{v}_{1}\}\mathbf{T} + \{R_{2} \mid \mathbf{v}_{2}\}\mathbf{T} + \{\overline{R}_{2} \mid \mathbf{v}_{2}\}\mathbf{T} + \dots\\
&\quad + \{R_{h} \mid \mathbf{v}_{h}\}\mathbf{T} + \{\overline{R}_{h} \mid \mathbf{v}_{h}\}\mathbf{T},
\end{array}
\tag{6.3.2}
\end{equation*}
其中 $R_{i}$ 与 $\overline{R}_{i}$ 是双点群中对应于点群元素 $R_{i}$（$R_{1}, R_{2}, \dots, R_{h}$ 中）的元素，而 $\mathbf{T}$ 是布拉维格子的平移群。

\begin{figure}[H]
\centering
\includegraphics[width=0.86\textwidth]{c6_63_p467.jpg}
\caption*{6.3 节的表（\textbf{表 6.7}--6.10 等，原书第 454 页起扫描）}
\end{figure}
\begin{figure}[H]
\centering
\includegraphics[width=0.86\textwidth]{c6_63_p468.jpg}
\caption*{原书第 455 页}
\end{figure}
\begin{figure}[H]
\centering
\includegraphics[width=0.86\textwidth]{c6_63_p469.jpg}
\caption*{原书第 456 页}
\end{figure}
\begin{figure}[H]
\centering
\includegraphics[width=0.86\textwidth]{c6_63_p470.jpg}
\caption*{原书第 457 页}
\end{figure}
\begin{figure}[H]
\centering
\includegraphics[width=0.86\textwidth]{c6_63_p471.jpg}
\caption*{原书第 458 页}
\end{figure}
\begin{figure}[H]
\centering
\includegraphics[width=0.86\textwidth]{c6_63_p472.jpg}
\caption*{原书第 459 页}
\end{figure}
\begin{figure}[H]
\centering
\includegraphics[width=0.86\textwidth]{c6_63_p473.jpg}
\caption*{原书第 460 页}
\end{figure}
\begin{figure}[H]
\centering
\includegraphics[width=0.86\textwidth]{c6_63_p474.jpg}
\caption*{原书第 461 页}
\end{figure}
\begin{figure}[H]
\centering
\includegraphics[width=0.86\textwidth]{c6_63_p475.jpg}
\caption*{原书第 462 页}
\end{figure}
\begin{figure}[H]
\centering
\includegraphics[width=0.86\textwidth]{c6_63_p476.jpg}
\caption*{原书第 463 页}
\end{figure}
\begin{figure}[H]
\centering
\includegraphics[width=0.86\textwidth]{c6_63_p477.jpg}
\caption*{原书第 464 页}
\end{figure}
\begin{figure}[H]
\centering
\includegraphics[width=0.86\textwidth]{c6_63_p478.jpg}
\caption*{原书第 465 页}
\end{figure}
\begin{figure}[H]
\centering
\includegraphics[width=0.86\textwidth]{c6_63_p479.jpg}
\caption*{原书第 466 页}
\end{figure}

% 6.4 推导双空间群表示的一个例子
\section{推导双空间群表示的一个例子}

在 5.4 节中我们演示了如何由表 5.7 得到空间群 $F\bar{4}3c\ (T^{5}_{d})$ 的单值表示，以此说明该表的用法。现在我们说明如何用同一张表推导同一空间群 $F\bar{4}3c\ (T^{5}_{d})$ 的双值表示。

空间群 $F\bar{4}3c\ (T^{5}_{d})$ 的布里渊区示于图 3.14 与图 5.2。由表 3.6，各对称点与对称线的 $\mathbf{k}$ 矢量以 $\mathbf{g}_{1}, \mathbf{g}_{2}, \mathbf{g}_{3}$ 为单位：
\begin{center}
\begin{tabular}{lc}
点或线 & $\mathbf{k}$\\
\midrule
$\Gamma$ & $(0, 0, 0)$\\
$X$ & $(\frac{1}{2}, 0, \frac{1}{2})$\\
$L$ & $(\frac{1}{2}, \frac{1}{2}, \frac{1}{2})$\\
$W$ & $(\frac{1}{2}, \frac{1}{4}, \frac{3}{4})$\\
$\Delta$ & $(\alpha, 0, \alpha)$\\
$\Lambda$ & $(\alpha, \alpha, \alpha)$\\
$\Sigma$ & $(\alpha, \alpha, 2\alpha)$\\
$S$ & $(\frac{1}{2}+\alpha, 2\alpha, \frac{1}{2}+\alpha)$\\
$Z$ & $(\frac{1}{2}, \alpha, \frac{1}{2}+\alpha)$\\
$Q$ & $(\frac{1}{2}, \frac{1}{2}-\alpha, \frac{1}{2}+\alpha)$
\end{tabular}
\end{center}

我们仍逐一考虑各对称点。

\noindent\textbf{$\Gamma$：$\mathbf{k} = (0, 0, 0)$。}

${}^{H}\mathbf{G}^{\Gamma}$ 是 $\mathbf{G}^{7}_{24}$，即点群 $\bar{4}3m\ (T_{d})$（见 5.4 节），因此 ${}^{H}\mathbf{G}^{\Gamma} \uparrow$ 所需的是 $\bar{4}3m\ (T_{d})$ 的双群，见表 6.5。该群也可辨认为抽象群 $\mathbf{G}^{10}_{48}$，其生成元可辨认为
\begin{equation*}
P = \{S^{-}_{4x} \mid {\textstyle\frac{1}{2}}{\textstyle\frac{1}{2}}{\textstyle\frac{1}{2}}\}, \quad Q = \{\overline{C}^{-}_{31} \mid 000\}, \quad R = \{\sigma_{db} \mid {\textstyle\frac{1}{2}}{\textstyle\frac{1}{2}}{\textstyle\frac{1}{2}}\}.
\end{equation*}
该群的八个类从而为：$C_{1}$：$\{E \mid 000\}$；$C_{2}$：$\{\overline{E} \mid 000\}$；$C_{3}$：$\{C_{2x} \mid 000\}, \{C_{2y} \mid 000\}, \{C_{2z} \mid 000\}$ 及其带杠版本；$C_{4}$、$C_{5}$、$C_{6}$、$C_{7}$、$C_{8}$ 同理按 $\bar{4}3m\ (T_{d})$ 的八个类加倍得到。于是 ${}^{H}\mathbf{G}^{\Gamma}$ 的特征标表可由表 5.1 的 $\mathbf{G}^{10}_{48}$ 条目、结合表 6.5 的标记得到；表示分记为 $\Gamma_{1}, \dots, \Gamma_{6}$（单值，$A_{1}, A_{2}, E, T_{1}, T_{2}$）与 $\Gamma_{7}$（双值，二维）等。由表 6.14 的“a”条目可查得其标记。各表示的实性（第一/二/三种）同样在表 6.13 中给出，从而可以判断计及时间反演对称后是否出现额外简并：在 $\Gamma$ 点，双值表示属于第一种，没有额外简并。

对其余各点 $X, L, W$ 与各对称线 $\Delta, \Lambda, \Sigma, S, Z, Q$，做法完全相同：由表 6.13 辨认 ${}^{H}\mathbf{G}^{\mathbf{k}}$（对称点）或 $\overline{\mathbf{G}}^{\mathbf{k}*}$（对称线）对应的抽象群、生成元及其在双群中的推广，再从表 6.1、6.5 与表 6.14 得到特征标、矩阵代表与标记。原书第 473--479 页按此逐一给出，此处不再重复；所得的双值小表示及其种类的分类可直接由表 6.13 与 6.14 检索。
